% !PRINTED 287
\weberpage{287}
Now if, at one of the points of intersection of the curves $\varphi$, $\psi$,
the positive direction of advance on $\varphi$ leads into the interior
of $\psi$, then $d\psi$, and hence also the functional determinant
$[\varphi, \psi]$, is negative, and in the opposite case positive (Fig.~7).
We call a point of intersection of $\varphi$ and $\psi$ a
\gesperrt{point of exit} $A(\varphi, \psi)$ or a point of entry
$E(\varphi,\psi)$, according as

\weberfigureimage{weber-b1-fig06-07.png}{0.94}

the positive direction of advance of $\varphi$ leads from the interior of $\psi$
to the exterior or from the exterior to the interior, and thus we have the theorem:
\weq{4}{
\begin{array}{lll}
\text{at an } A(\varphi, \psi) & \text{is} & [\varphi, \psi] > 0 \\
\text{,,}\ \ \text{,,} & E(\varphi, \psi) & \text{,,}\ \ [\varphi, \psi] < 0.
\end{array}
}
The number of the $A(\varphi, \psi)$ is just as great as that of the $E(\varphi, \psi)$;
and if we interchange $\varphi$ with $\psi$, the $A(\varphi, \psi)$, $E(\varphi, \psi)$
pass into $E(\psi, \varphi)$ and $A(\psi, \varphi)$.

The two functions $\varphi$, $\psi$ determine in a unique manner
a region of the plane, which is characterised by the fact that in it
the product $\varphi\psi$ is negative. We shall call it the \gesperrt{inner region}
$(\varphi, \psi)$ or $(\psi, \varphi)$. In our Fig.~7 it is the shaded
area.

\begin{center}
\S.\,94.

\gesperrt{Characteristic of a system of three functions.}
\end{center}

We now adjoin to the two functions $\varphi$, $\psi$ a third,
$f(x, y)$. The curve represented by the equation $f = 0$,
or the curve $f$, shall likewise be closed in the finite part of the plane
%CONT>>
